% Independent targeted re-proof, 2026-11-19, agent period_triangle.
\section*{Primary-source re-proof of the first assertion of LM6 A.20}
The assertion proved here is the all-prime procyclic normalizer recognition statement for $E_0:y^2=x^3-x$ over $K=\mathbf Q(i)$. The separate integral claim about an actual class-four lift at $\ell=5$ is not used and is not certified by this proof.

Let $U=E_0\setminus\{O\}$, $F=\pi_1(U_{\overline K})$, $P_\ell=F^{(\ell)}$, and
\[
 E_\ell=\pi_1(U)/\ker(F\to P_\ell).
\]
Let $A\subset P_\ell$ be a nontrivial closed procyclic subgroup and let $N=N_{E_\ell}(A)$ have open image in $G_K$. Then $A$ is contained in a cusp inertia subgroup of $P_\ell$.

\paragraph{Primary sources actually opened.}
Ishii, arXiv:2312.04175v2, 4 September 2024, standing assumptions on printed pp.4--5, Theorem1.3 on p.3, and Corollary4.10 on p.21. The journal landing page for the 2025 publication was opened but supplies only a subscription preview; the full text actually used is the author's arXiv version. Nakamura, \emph{On exterior Galois representations associated with open elliptic curves}, J. Math. Sci. Univ. Tokyo2 (1995), 197--231: (4.3)--(4.4), printed p.214; Proposition4.7 and (4.8), p.215; exact sequences (4.9.1)--(4.9.2), p.216; Lemma4.10, p.217; Corollary4.12, p.221. Nakamura, \emph{Galois rigidity of pure sphere braid groups and profinite calculus}, same journal1 (1994), 71--136, Theorem2.1.1 on p.94, for the final full pro-$\ell$ inertia recognition. The last source's definition includes the cyclotomic action as an explicit premise; that premise is proved below.

\paragraph{1. Actual finite-cover homology detects a generator.}
The free pro-$\ell$ group $P_\ell$ is torsion-free, so $A\simeq\mathbf Z_\ell$; fix a generator $u$. Choose a finite $\ell$-group quotient $P_\ell\to Q$ in which $u$ has nonidentity image $\bar u$. For
\[
 H=\operatorname{preimage}(\langle\bar u\rangle)\subset P_\ell,
\]
the composite $H\to\langle\bar u\rangle\to\mathbf Z/\ell$ sends $u$ to a nonzero element. Thus $u$ has nonzero image in the torsion-free abelianization $H^{\rm ab}$.

The inverse image of $H$ in $F$ defines an actual finite etale geometric cover $Y\to U_{\overline K}$. It descends to a finite extension of $K$. Moreover $\pi_1(Y_{\overline K})^{(\ell)}\simeq H$. To check this exact identification, any finite $\ell$-group quotient of $\pi_1(Y_{\overline K})$ has normal closure over $U_{\overline K}$ with an $\ell$-group as Galois group: take the Galois closure of $Y$, which is already an $\ell$-cover, and the compositum of its finitely many conjugate $\ell$-covers. The resulting group is a subgroup of an iterated wreath product of finite $\ell$-groups. Hence all these quotients factor through $P_\ell$, proving the asserted maximal-pro-$\ell$ identification.

Take an arithmetic covering subgroup for a descent of $Y$ and intersect it with $N$. Its arithmetic image is still open, since that covering subgroup is open. Its geometric kernel acts inner on $H$ and hence trivially on $H^{\rm ab}$. Therefore the nonzero line spanned by $u$ in
\[
 H^{\rm ab}\otimes\mathbf Q_\ell
 \simeq H_1(Y_{\overline K},\mathbf Q_\ell)
\]
is stable for the ordinary Galois action of a finite number-field extension. This construction uses neither a section of $N$ nor a splitting of an arithmetic extension. The homology exact sequence for the smooth compactification $\overline Y$ has boundary term of weight two and proper term of weight one. Thus the normalizer character on this line has Frobenius weight one or two.

\paragraph{2. The first nonzero term has degree one or two.}
Use the completed free Lie algebra over $\mathbf Q_\ell$ on $V=H_1(E_{0,\overline K},\mathbf Q_\ell)$. The pro-$\ell$ Magnus embedding is injective, and its logarithm gives a nonzero series $U_u=\log u$ with first nonzero homogeneous term $u_d\in L_d(V)$. For clarity, the Magnus embedding here comes from the canonical completed-group-algebra identification
\[
 \mathbf Z_\ell[[P_\ell]]\simeq\mathbf Z_\ell\langle\!\langle X,Y\rangle\!\rangle,
 \qquad x_0\mapsto1+X,\quad y_0\mapsto1+Y.
\]
The group is separated by its finite quotients, so its inclusion in the units of its completed group algebra is faithful. After the formal substitutions $x=\log(1+X)$, $y=\log(1+Y)$, the logarithm of a group-like element is primitive, hence lies in the completed free Lie algebra. No integral division by two or by $\ell$ is used to infer a group identity.

Conjugation by a geometric element acts trivially on the first homogeneous term. If $n\in N$ acts by $u\mapsto u^{\theta(n)}$, then
\[
 p(n)(u_d)=\theta(n)u_d.
\]
The character therefore descends to $p(N)$, and at a good arithmetic Frobenius of residue cardinality $q$, its complex absolute value is $q^{d/2}$, since $V$ has weight one. Comparison with the actual detector in Step1 gives $d=1$ or $d=2$.

\paragraph{3. Degree one at inert primes and at two.}
The curve has trace $-2$ over $\mathbf F_5$, so a good Frobenius at a prime over5 has eigenvalues
\[
 \alpha=-1+2i,\qquad\bar\alpha=-1-2i.
\]
Their ratio is not a root of unity: it lies in $\mathbf Q(i)$ and is none of its four roots of unity $1,-1,i,-i$. A positive power of this Frobenius lies in any prescribed open subgroup of $G_K$. When $\ell=2$ or $\ell$ is inert in $K$, $i\notin\mathbf Q_\ell$, and no positive power of either eigenvalue belongs to $\mathbf Q_\ell$, since equality with its conjugate would make $(\alpha/\bar\alpha)^n=1$. Thus the characteristic polynomial of every such power is irreducible over $\mathbf Q_\ell$, and $V$ has no one-dimensional subspace stable under an open arithmetic subgroup. This excludes $d=1$ at those primes.

\paragraph{4. The split-prime source input and its exact degree.}
Now let $\ell$ split in $K$, hence $\ell\ge5$. The CM eigenspaces give a decomposition $V=\mathbf Q_\ell x\oplus\mathbf Q_\ell y$, with $[i]$ acting by $i$ and $-i$. Every line stable under an open arithmetic subgroup is one of these two lines. Indeed the two CM characters remain distinct after finite extension; the preceding Frobenius5 argument proves this for $\ell\ne5$, and at $\ell=5$ use Frobenius13, whose eigenvalues are $3\pm2i$ and whose ratio is again not a root of unity.

Put $K_\infty=K(E_0[\ell^\infty])$. Nakamura's (4.3)--(4.4) gives a canonical representative of every outer Galois action from $G_{K_\infty}$ which fixes the boundary and is the identity through derivation degree three:
\[
 G_{K_\infty}\longrightarrow\Gamma^*_{1,1}(4).
\]
This particular representative matters: reduction to a metabelian group without retaining the boundary-fixing representative would lose the following information. Nakamura's (4.7)--(4.9) defines the comparison
\[
 \operatorname{gr}_4\Gamma^*_{1,1}\longrightarrow
 \operatorname{gr}_4\Psi^*\simeq\operatorname{Sym}^2V\otimes\det(V).
\]
It is surjective by Lemma4.10 for odd $\ell$. The exact sequence (4.9.2) gives
\[
 \dim\operatorname{gr}_4\Gamma^*_{1,1}
 =2\dim L_5(V)-\dim L_6(V)=2\cdot6-9=3.
\]
The target also has dimension3, so the displayed map is an isomorphism after tensoring with $\mathbf Q_\ell$.

For $[i]$, the three weight spaces in its target have characters $(3,1),(2,2),(1,3)$ and eigenvalues $-1,1,-1$. The arithmetic outer action commutes with the $K$-defined endomorphism $[i]$, so its degree-four image lies in the middle line. Ishii's Corollary4.10 gives a nonzero character $\kappa_{(2,2)}$ on $G_{K_\infty}$. All standing assumptions hold: $K=\mathbf Q(i)$ has class number one, the full CM is defined over $K$, and $\ell\ge5$ splits. The curve has discriminant64, so it has good reduction at $\ell$. Its Hasse invariant is the coefficient of $x^{\ell-1}$ in $(x^3-x)^{(\ell-1)/2}$, namely $(-1)^{(\ell-1)/4}\binom{(\ell-1)/2}{(\ell-1)/4}$ for $\ell\equiv1\pmod4$; this is nonzero modulo $\ell$, verifying ordinary reduction directly. His Theorem1.3, equivalently Nakamura's Corollary4.12 with the shifted index convention, places this character in the coefficient of $U_1U_2$ in the quadratic part of the comparison power series, multiplied by the nonzero factor $(1-\ell^2)^{-1}$. Thus the middle degree-four outer action is nonzero.

A nonzero continuous homomorphism to the torsion-free additive group $\mathbf Z_\ell$ remains nonzero on every open subgroup: otherwise its image would be finite. In particular there exists
\[
 \sigma\in G_{K_\infty}\cap p(N)
\]
with nonzero middle degree-four action. This supplies the required arithmetic element without choosing a section of $N\to p(N)$.

\paragraph{5. Independent calculation of the middle derivation.}
Let $\omega=[x,y]$, $F_0(x)=y$, $F_0(y)=0$, and define
\[
 \epsilon_4(x)=\operatorname{ad}_x^4y,\quad
 \epsilon_4(y)=[y,\operatorname{ad}_x^3y]-[[x,y],\operatorname{ad}_x^2y],
 \qquad T_4=[F_0,\epsilon_4].
\]
Direct expansion gives $T_4(\omega)=0$ and, in standard Lyndon brackets,
\[
 T_4(x)=-2b_{xxxyy}-b_{xxyxy},\qquad
 T_4(y)=2b_{xxyyy}+3b_{xyxyy}.
\]
The three special derivations $\epsilon_4,T_4,[F_0,T_4]$ form a basis: the kernel of $(D(x),D(y))\mapsto[D(x),y]+[x,D(y)]$ has dimension3, and these three pairs are linearly independent. In the metabelian Lie quotient, with $L'/L''=\mathbf Q_\ell[X,Y]\omega$, their pairs have the form
\[
 (D(x),D(y))=(X\beta\omega,Y\beta\omega),\qquad
 \beta=X^2,\ 2XY,\ 2Y^2
\]
respectively. Hence their comparison determinant is4 and the middle line is exactly $\mathbf Q_\ell T_4$. The sign of $\omega$ relative to Nakamura's inverse-commutator boundary changes only a nonzero scalar and is immaterial to nonvanishing.

Most importantly,
\[
 T_4(x)\notin[x,L_4],\qquad T_4(y)\notin[y,L_4].
\]
In the relevant bidegree, each bracket image has dimension1, and adjoining the indicated vector raises the rank to2. The independently written exact script \texttt{check\_CM\_degree\_four.py} verifies the kernel dimension, all three special derivations, the three metabelian polynomials, and both rank obstructions. It also independently counts the reductions at5 and13. These computations are finite algebra certificates; the preceding primary-source result supplies the arithmetic nonvanishing.

\paragraph{6. The full conjugacy equation, including lower inner terms.}
Suppose $d=1$ and, after interchanging the two CM lines if needed, write
\[
 U_u=\log u=a x+\text{terms of degree at least2},\qquad a\in\mathbf Q_\ell^\times.
\]
No primitivity of the original abelian image is assumed. Choose the $\sigma$ in Step4 and a lift $n\in N$. Since $\sigma$ acts trivially on $V$ and $a\ne0$, the normalizer exponent is $\theta(n)=1$. Let $\widetilde\sigma$ denote Nakamura's boundary-fixing representative. Its first nonidentity derivation is $cT_4$ with $c\ne0$, whereas the representative furnished by $n$ differs from it by an inner automorphism. Thus, in the rational completed free Lie algebra,
\[
 \widetilde\sigma(U_u)=e^{\operatorname{ad} W}U_u
\]
for some $W$.

Comparing degree2 first gives $[W_1,x]=0$, hence $W_1=b x$. Replace the conjugating element $\exp W$ on the right by
\[
 \exp W\,\exp\bigl(-(b/a)U_u\bigr).
\]
This preserves its conjugation action on $U_u$ and removes its degree-one term. Comparing degrees3 and4 now forces the degree-two and degree-three terms to vanish. The fact used here is $C_L(x)=\mathbf Q_\ell x$: in the free associative algebra a homogeneous element commuting with $x$ is a scalar multiple of a power of $x$, and among these only the degree-one term is a nonzero primitive Lie element in characteristic zero.

Consequently the degree-five equation is
\[
 acT_4(x)=[W_4,a x],
\]
contradicting Step5. The same argument with $y$ uses the second rank obstruction. This excludes degree one for every split prime, and completes the part of the prior argument that had left its inner correction implicit.

\paragraph{7. Degree two gives the actual cyclotomic action and actual inertia.}
Thus $d=2$. Since $L_2(V)=\mathbf Q_\ell\omega$, write $u_2=a\omega$ with $a\ne0$. For every $n\in N$, comparison on this term gives
\[
 a\theta(n)\omega=a\chi_\ell(p(n))\omega.
\]
Canceling the nonzero scalar yields $\theta=\chi_\ell\circ p$ on the actual normalizer. Nakamura1994, Theorem2.1.1, applied to the full class of finite $\ell$-groups, now puts $u$ and $A$ in an actual cusp inertia subgroup of $P_\ell$. The conclusion concerns the pro-$\ell$ group itself; the rational Lie algebra was used only to rule out the degree-one case.

\paragraph{Boundary of this audit.}
This independently certifies the first recognition assertion of LM6 A.20 and therefore the compositional use of that assertion in the explicit quartic cover. It does not certify the separate class-four integral lifting statement at5, does not infer an integral normalization such as $\kappa/48$, and does not infer that all remaining full-profinite A.22 subgroups are detected in one of the displayed covers.
